\section{Uniqueness of the actual canonical first-meeting prefix}

All prefixes in this section refer to the same original marked record and
the same original good open set.  Canonical histories may contain different
numbers of actual empty steps.  Comparisons concern their genuine maps over
the original base, rather than equality of their stopping indices.

\begin{theorem}[Actual empty compression preserves the avoided original open]
\label{hir:original-empty-compression-preserves-first-open}
\uses{hir:original-finite-family-empty-free-canonical-compression,hir:original-actual-empty-family-canonical-deletion,hir:tower-isomorphism}
\lean{Hironaka.HironakaOriginalFirstMeetingPrefixUniqueness.empty_free_compression_preserves_avoidance}
\leanok
\source{wlodarczyk05-hironaka,bierstone-milman08-functoriality}
Let $n>0$, let $M$ belong to the original invariant domain of bound $k$,
and let $T$ be an actual canonical history for the unchanged recursive
family of bound $k$.  If all its centres avoid an original open $O$, one
can internally remove every actual empty step to obtain a canonical
history $Q$ whose actual centres are nonempty and still avoid $O$.
There are a biholomorphism $\chi$ from the endpoint of $Q$ to that of $T$
over the original base and one order embedding $F$ such that the whole
cosupport and every current invariant value on it are transported by
$\chi$ and $F$.  In particular, the next whole maximum of $T$ meets its
inverse image of $O$ if and only if that of $Q$ does.
\end{theorem}
\begin{proof}
Induct on the actual finite length.  If every centre is nonempty, retain
the same tower, identity biholomorphism and identity order embedding.
Otherwise split off a genuine empty step and its whole suffix.  The
actual empty blowup is biholomorphic to its base; transport the entire
suffix through this biholomorphism.  This identifies the original
history with the transported history containing an explicit empty
identity step.  Centre avoidance is invariant under this genuine
stagewise tower isomorphism: the centres and their projection inverse
images correspond, and the stage maps are bijective.  Removing the
empty identity step preserves avoidance, since it contributes no centre
and does not change any subsequent projection to the original base.
The actual empty-deletion theorem supplies canonicality of the shorter
history and its whole cosupport and value dictionary.  Apply induction
to this strictly shorter actual history and compose its biholomorphism
and order embedding with those from the deletion.  This terminates and
gives the claimed whole dictionaries and avoided open.  A bijection of
cosupports and an order embedding preserve and reflect the defining
inequalities against every cosupport point, so they transport the whole
maximum locus.  The over-base projection identity then preserves its
encounter with $O$ in both directions.
\end{proof}

\begin{theorem}[Avoidance and the next whole maximum determine the nonempty length]
\label{hir:original-first-meeting-nonempty-length}
\uses{hir:canonical-resolution,hir:tower-isomorphism}
\lean{Hironaka.HironakaOriginalFirstMeetingPrefixUniqueness.first_meeting_nonempty_lengths_eq}
\leanok
\source{wlodarczyk05-hironaka,bierstone-milman08-functoriality}
For an original resolution invariant and an original record in its
domain, let $T_1,T_2$ be canonical segments with nonempty actual centres.
If both avoid the same original open $O$, and each next whole maximal
locus meets the inverse image of $O$, then $T_1,T_2$ have equal lengths.
Consequently they are genuinely isomorphic towers over the original base.
\end{theorem}
\begin{proof}
Suppose $T_1$ were longer.  Split it after exactly the number of steps of
$T_2$.  Equal-length uniqueness for nonempty canonical segments gives a
genuine stagewise tower isomorphism between $T_2$ and this prefix.  The
whole cosupport and invariant dictionary transports the next maximal
locus, so it meets the inverse image of $O$ at this prefix.  Its remaining
tail is nonempty.  Its first actual centre is exactly that whole maximum,
because the segment has nonempty canonical centres.  This centre meets
the inverse image of $O$, contradicting avoidance of the whole $T_1$.
The same argument with the roles reversed excludes $T_2$ being longer.
Thus their lengths agree, and equal-length canonical uniqueness supplies
the claimed tower isomorphism.
\end{proof}

\begin{theorem}[The original canonical first-meeting map is intrinsically unique]
\label{hir:original-first-meeting-prefix-map-unique}
\uses{hir:original-empty-compression-preserves-first-open,hir:original-first-meeting-nonempty-length,hir:original-actual-whole-projection-biholomorphic-map-uniqueness}
\lean{Hironaka.HironakaOriginalFirstMeetingPrefixUniqueness.first_meeting_prefixes_existsUnique_biholomorphism}
\leanok
\source{wlodarczyk05-hironaka,bierstone-milman08-functoriality}
For the original ambient canonical invariant, let $T_1,T_2$ be any two
actual finite canonical prefixes of the same original marked record in
its invariant domain.  Suppose they avoid the same original open $O$ and
their next whole maximal loci meet the inverse images of $O$.  There is
a unique biholomorphism $\eta$ between their original endpoints over the
unchanged original base.  No terminality, common actual length, supplied
compression, comparison map, or equality of stopping indices is a premise.
\end{theorem}
\begin{proof}
In positive dimension, compress both actual histories by the first
theorem, retaining avoidance, the next-maximum encounter, and the two
actual over-base endpoint biholomorphisms.  Their nonempty lengths agree
by the preceding theorem.  Compare these compressed canonical segments
by their actual tower isomorphism, then compose with the inverse and
forward compression biholomorphisms.  Their whole projection formulas
give the asserted equality over the unchanged original base.  Any two
such biholomorphisms agree by the actual finite-tower projection
uniqueness theorem, which uses the genuine successive blowup universal
properties.  In ambient dimension zero the original invariant-domain
predicate forces every transformed cosupport of a canonical history to
be empty; the stipulated next maximal point cannot exist.  This covers
all dimensions without imposing a false positive-dimension bound.
\end{proof}

\begin{remark}
This theorem applies to independently chosen embedded first-meeting
prefixes on the same original marked record and good open, including
different actual empty-step counts.  Together with the actual smooth
pullback comparison it concerns the same embedded stopped projection.
Relabeling a singular source's ideal as its separately reduced vanishing
ideal would additionally require a whole singular vanishing-ideal
pullback theorem.  No such relabeling or arbitrary smooth-source stopping
index law is asserted here.
\end{remark}
